\begingroup
\def\C{\mathbb{C}}
\def\R{\mathbb{R}}
\def\Tr{\operatorname{Tr}}
\section{The first canonical window, its full phase product, and both arithmetic block budgets}



This chapter joins the original first-degree source calculation
(CV.11)--(CV.12) and (TVB.37)--(TVB.38) to the consecutive-window
proof of PR~25 at the exact head
\texttt{66fac5e7885a40cd4873}\allowbreak\texttt{e74b754907320f05b067}.
The complete original and proposed corrected sources accompany this
chapter. Their correction retains all zero generalized eigenvalues.
The source radius, original norm envelope and finite propagation
are proved below with their original domains. The full first-window
phase product has a single terminal contraction and two copies of
each interior contraction; its first step contains no preceding
quotient inverse.

\subsection{Original source and the actual first-degree step}

Retain $g=2\xi$, the complete fixed packet $h$, integer tensor degree $k\ge1$,
and the original density
\[
 w_h(t)=\frac{|(g/h)(1/2+it)|^2}{2\pi},\qquad
 m_{h,k}=w_h^{*k},\qquad S=c+iu,\quad c=k/2.
 \tag{CJ.1}
\]
Its source inner product is the integral against $m_{h,k}(u)du$;
the original monic orthogonal polynomials $p_n$ have squared norms
$\omega_n>0$, including $\omega_0=\mu_h^k$.
Use the original dagger-stable monic cyclic polynomial $\chi$ of
degree $q\ge1$, the full quotient $E=\C[S]/(\chi)$ and the
original remainder basis. Its columns, kernels and metrics are
\[
 \begin{gathered}
 b_n=[p_n]_\chi,\quad B_N=[b_0,\ldots,b_N],\quad
 O_N=\operatorname{diag}(\omega_0,\ldots,\omega_N),\quad
 \\
 K_N=B_NO_N^{-1}B_N^*,\quad G_N=K_N^{-1},\quad V_N=\det G_N
 \quad(N\ge q-1).
 \end{gathered}
 \tag{CJ.2}
\]
Surjectivity of $B_N$ and positivity of $O_N$ prove $K_N>0$.
The section $R_N=O_N^{-1}B_N^*G_N$ satisfies $B_NR_N=I$ and
$R_N^*O_NR_N=G_N$; its orthogonality to $\ker B_N$ proves
that it is the unique minimum lift. Source inclusion proves
$G_{N+1}\preceq G_N$ and $0<V_{N+1}\le V_N$.

Write the original expansion, retaining every coefficient,
\[
 \chi=p_q+\sum_{j=0}^{q-1}\gamma_jp_j,
 \qquad \gamma_j=\frac{\langle p_j,\chi\rangle}{\omega_j},
 \qquad \nu_0=\|\chi\|^2
       =\omega_q+\sum_{j<q}|\gamma_j|^2\omega_j.
 \tag{CJ.3}
\]
The first $q$ remainder columns form an invertible triangular
matrix, so $b_a^*G_{q-1}b_b=\delta_{ab}\omega_b$ for $a,b<q$.
Reduction of (CJ.3) gives
$b_q=-\sum_{j<q}\gamma_jb_j$ and hence
$b_q^*G_{q-1}b_q=\nu_0-\omega_q$.
The exact rank-one determinant formula therefore gives
\[
 \delta_q:=\frac{V_q}{V_{q-1}}
 =\left(1+\frac{b_q^*G_{q-1}b_q}{\omega_q}\right)^{-1}
 =\frac{\omega_q}{\nu_0}.
 \tag{CJ.4}
\]
The rank-one determinant identity follows by taking a basis whose
first vector is the update vector $u$ when it is nonzero.
In that basis $uv^*$ has entries only in its first row;
$I+uv^*$ is therefore upper triangular with diagonal entries
$1+v^*u,1,\ldots,1$. Its determinant is $1+v^*u$.
The zero update gives the same identity directly.

Let $\phi_N$ denote the retained real phase of the original
theta deformation, with the sign convention of (TVB.26):
$b_N^*G_Nb_{N+1}=i\omega_N\phi_N$.
Its reality follows in the literal chart $S=c+iu$ as follows.
Write $Q_n(u)=i^{-n}p_n(c+iu)$ and
$\widehat\chi(u)=i^{-q}\chi(c+iu)$. Dagger stability gives
$\overline\chi(k-S)=(-1)^q\chi(S)$, so $\widehat\chi$
is real monic. The source measure is real and positive; the
monic orthogonality equations therefore give real $Q_n$.
If $\mathsf E$ is the exact coefficient map
$P(S)\mapsto P(c+iu)$, then
$\mathsf E b_n=i^n[Q_n]_{\widehat\chi}$, while the transported
kernel and metric are $\mathsf E K_N\mathsf E^*$ and
$\mathsf E^{-*}G_N\mathsf E^{-1}$.
The latter has real entries because the phases in each kernel
summand occur together with their conjugates. Thus the original
cross pairing is $i$ times a real number, with its original
factor $\omega_N$ retained. This proves the stated reality
and specifies the exact chart map rather than changing the metric.
At the first degree that identity reads
$\gamma_{q-1}=-i\phi_{q-1}$. The original two-column Hermitian
control formula (CV.4), before specialization, is
\[
 \epsilon_N^2=
 \frac{a_N d_N-|z_N|^2}{\omega_N^2},\qquad
 a_N=b_N^*G_Nb_N,\quad z_N=b_N^*G_Nb_{N+1},\quad
 d_N=b_{N+1}^*G_Nb_{N+1}.
 \tag{CJ.5}
\]
Here $z_N$ is purely imaginary in the stated dagger-stable
source. The original boundary-control identity on $(E,G_N)$ is
\[
 \mathsf H_N=G_N^{-1}(A^*G_N+G_NA-kG_N)
 =\frac{b_{N+1}b_N^*G_N+b_Nb_{N+1}^*G_N}{\omega_N},
 \qquad A=M_S.
\]
For completeness this identity follows directly from the original
monic recurrence. Self-adjointness of multiplication by $u$ in
the integral pairing and orthogonality give
$uQ_n=Q_{n+1}+\beta_nQ_n+\alpha_nQ_{n-1}$, where
$\beta_n\in\R$ and $\alpha_n=\omega_n/\omega_{n-1}$ for
$n\ge1$; the lower term is zero at $n=0$.
Indeed all coefficients below degree $n-1$ vanish by pairing
with $uQ_j$ of degree at most $n-1$, and pairing with
$Q_{n-1}$ gives the displayed norm ratio. In the original
$S$ coordinates this is
$Sp_n=p_{n+1}+(c+i\beta_n)p_n-\alpha_np_{n-1}$.
Reduce it modulo $\chi$ and insert it into
$AK_N+K_NA^*-kK_N$. The diagonal coefficients are
$c+i\beta_n+c-i\beta_n-k=0$. For each internal adjacent
pair the coefficient is
$\omega_n^{-1}-\alpha_{n+1}\omega_{n+1}^{-1}=0$.
The remaining two terms are
$(b_{N+1}b_N^*+b_Nb_{N+1}^*)/\omega_N$.
Multiplication on the right by the original $G_N$ proves
the written boundary-control formula with its signs.
It has rank at most two, is self-adjoint in $G_N$, and has
trace $(z_N+\overline z_N)/\omega_N=0$. Squaring its two
summands and taking the trace gives
$[z_N^2+\overline z_N^2+2a_Nd_N]/\omega_N^2
=2(a_Nd_N-|z_N|^2)/\omega_N^2$.
For $q\ge2$ its full spectrum is $\epsilon_N,-\epsilon_N$
and $q-2$ zeros, with $\epsilon_N\ge0$; these first two
entries may themselves be zero. For $q=1$ trace zero makes
the entire control zero. This proves (CJ.5), including the cross-pairing.
At $N=q-1$ the three entries are respectively
$\omega_{q-1},-\gamma_{q-1}\omega_{q-1},\nu_0-\omega_q$.
Substitution proves
\[
 \epsilon_{q-1}^2+\phi_{q-1}^2
 =\frac{\nu_0-\omega_q}{\omega_{q-1}}.
 \tag{CJ.6}
\]
For $q=1$, the lower projection in (CV.12) is the zero source
and $\epsilon_{q-1}=0$; the retained phase identity and (CJ.6)
still hold. No $G_{q-2}$ occurs in this calculation.

Define $\Lambda_N=\omega_N/V_N$. Combining (CJ.4)--(CJ.6)
gives the exact first step
\[
 (\epsilon_{q-1}^2+\phi_{q-1}^2)\Lambda_{q-1}
      =\Lambda_q(1-\delta_q).
 \tag{CJ.7}
\]
Indeed its left side is $(\nu_0-\omega_q)/V_{q-1}$, and
$\nu_0/V_{q-1}=\omega_q/V_q$ by (CJ.4). This proves the
identity with its single contraction and full original phase.

\subsection{The complete first-window product and both finite budgets}

At each regular degree $N\ge q$, the proved original radius
identity (EP.11), equivalently (TVB.27), is
\[
 (\epsilon_N^2+\phi_N^2)\Lambda_N
 =\Lambda_{N+1}(1-\delta_N)(1-\delta_{N+1}),
 \qquad \delta_N=V_N/V_{N-1}.
 \tag{CJ.8}
\]
Every $\delta_N$ lies in $(0,1]$. Multiply (CJ.7) by the
consecutive identities (CJ.8) at $N=q,\ldots,j-1$.
All $\Lambda_N$ are positive, so canceling their identical
intermediate factors is valid. It yields, for every $j\ge q$,
\[
 \boxed{
 \prod_{N=q-1}^{j-1}(\epsilon_N^2+\phi_N^2)
 =\frac{\omega_jV_{q-1}}{\omega_{q-1}V_j}
       (1-\delta_j)\prod_{N=q}^{j-1}(1-\delta_N)^2.}
 \tag{CJ.9}
\]
For $j=q$ the interior product is empty and equal to one,
so this is precisely the first step. For $j>q$, each
$\delta_N$, $q\le N<j$, occurs once from the step ending at
$N$ and once from the step beginning there. The terminal
$\delta_j$ occurs once. This proves every exponent in (CJ.9).
For a regular initial degree $i\ge q$, the identical multiplication
instead gives
\[
 \prod_{N=i}^{j-1}(\epsilon_N^2+\phi_N^2)
 =\frac{\omega_jV_i}{\omega_iV_j}
       (1-\delta_i)(1-\delta_j)
       \prod_{N=i+1}^{j-1}(1-\delta_N)^2,
 \qquad j>i.
 \tag{CJ.10}
\]
These two expressions explicitly relate the first and regular
windows to the same source, phase and quotient volumes.

For the original off-line quartet let its common multiplicity
be $m\ge1$, its real displacement be $\delta>0$ and its
positive imaginary coordinate be $\gamma>0$. The complete
tensor calculation gives, with no omitted repeated root,
\[
 q=[1+k(m-1)](k+1)^2,\qquad
 L=2\delta[1+k(m-1)](k+1)
              \left\lfloor\frac{(k+1)^2}{4}\right\rfloor,
 \qquad L/q\ge\delta k/2.
 \tag{CJ.11}
\]
The original Hermitian compression proof gives
$0<L\le\epsilon_N$ at every admitted degree: on the full
sum of generalized eigenspaces with $\Re\lambda>k/2$, its
trace is $L$. An orthogonal projection onto that subspace in
the original metric has trace pairing at most $\epsilon_N$
with a Hermitian matrix of eigenvalues $\epsilon_N,-\epsilon_N$
and zeros. This is the same source metric as (CJ.2).
Thus each contraction occurring in (CJ.9)--(CJ.10) is strictly
between zero and one. A value one for any $\delta_N$ would
make a positive radius product zero.

For either original length-$q$ block $(i,j)=(q-1,2q-1)$ or
$(q,2q)$, taking logarithms of its exact product now gives
\[
 \log(V_i/V_j)
 =\sum_{N=i}^{j-1}\log(\epsilon_N^2+\phi_N^2)
       -\log(\omega_j/\omega_i)+\mathcal P_{i,j},
 \tag{CJ.12}
\]
where the entire contraction penalty is
\[
 \begin{split}
 \mathcal P_{q-1,2q-1}
 &=-\log(1-\delta_{2q-1})
       -2\sum_{N=q}^{2q-2}\log(1-\delta_N),\\
 \mathcal P_{q,2q}
 &=-\log(1-\delta_q)-\log(1-\delta_{2q})
       -2\sum_{N=q+1}^{2q-1}\log(1-\delta_N).
 \end{split}
 \tag{CJ.13}
\]
Both sums retain their literal first and last degrees and are
nonnegative. Since $\epsilon_N\ge L$, the exact finite bound is
\[
 \log(V_i/V_j)\ge2q\log L-\log(\omega_j/\omega_i)
   +\mathcal P_{i,j}
   +\sum_{N=i}^{j-1}\log(1+\phi_N^2/L^2).
 \tag{CJ.14}
\]
Every phase term remains; none is assigned the value zero.

\subsection{Actual comparable norm windows with complete constants}

The original arithmetic lower-envelope proof gives, for $k\ge3$,
\[
 m_{h,k}(u)\ge c_h\vartheta_h^{k-3}
 e^{-\alpha(k-3+|u|)}(k-2+|u|)^{-B},
 \quad\alpha=\pi/2,\quad B=42+2\deg h.
 \tag{CJ.15}
\]
Its complete analytic proof is retained in the arithmetic source
appendix and AU.1--AU.12. Fix $0<b<\pi/2$, and keep
$M_h(\pm b)=\int e^{\pm bt}w_h(t)dt$.
Set the explicit positive constants
\[
 \begin{gathered}
 t_h=\min(1,\vartheta_h),\quad c_0=\min(1,2c_h/3),\quad
 \ell_h=\tfrac12e^{-\alpha}\sqrt{c_0t_h}\,2^{-B/2},
 \\
 M=\max(1,M_h(b),M_h(-b)),\quad U_h=4\sqrt M/b.
 \end{gathered}
 \tag{CJ.16}
\]
For every integer $n\ge k\ge3$ the literal monic norms obey
\[
 (\ell_h n)^{2n}\le\omega_n\le(U_h n)^{2n}.
 \tag{CJ.17}
\]
Here is the full envelope calculation. The exact minimum squared
norm of a monic polynomial on $[-n,n]$ is
\[
 \frac{2n^{2n+1}}{2n+1}
       \left[\frac{2^n(n!)^2}{(2n)!}\right]^2
       \ge\frac23 n^{2n}4^{-n}.
\]
The last inequality uses $(2n)!/(n!)^2\le4^n$ and
$2n/(2n+1)\ge2/3$. The original chart's factor $i^n$ has
modulus one; its monic polynomial space is exactly this extremal
space through the inverse substitution. On $[-n,n]$, (CJ.15)
is at least
$c_h\vartheta_h^{k-3}e^{-\alpha(k-3+n)}(k-2+n)^{-B}$.
Use, in their stated directions,
\[
 \vartheta_h^{k-3}\ge t_h^n,\quad
 e^{-\alpha(k-3+n)}\ge e^{-2\alpha n},\quad
 (k-2+n)^{-B}\ge(2n)^{-B}\ge2^{-nB},\quad
 2c_h/3\ge c_0^n.
\]
The elementary $2n\le2^n$ holds at $n=1$ and follows by
induction from $2(n+1)\le4n$ for $n\ge1$.
Multiplication proves the lower bound (CJ.17).
For the upper bound use the monic trial polynomial $(S-c)^n$,
the nonnegative exponential series, and convolution to obtain
$\omega_n\le(2n)!b^{-2n}(M_h(b)^k+M_h(-b)^k)$.
This is at most $(2n)^{2n}b^{-2n}2M^n$, which is at most
$(4\sqrt M n/b)^{2n}$ since $2\le4^n$.
This proves both inequalities with the original mass retained.

For integers $n\ge k\ge3$ and $n/2\le r\le2n$, divide
the two instances of (CJ.17) and take the positive $2r$-th root:
\[
 \left(\frac{\omega_{n+r}}{\omega_n}\right)^{1/(2r)}
 \le n U_h^{1+n/r}\ell_h^{-n/r}
                  (1+r/n)^{1+n/r}
 \le C_h^{\rm win}n,
\]
\[
 C_h^{\rm win}=27\max(1,U_h)^3\max(1,\ell_h^{-1})^2.
 \tag{CJ.18}
\]
Indeed $1+n/r\le3$, $n/r\le2$, and $1+r/n\le3$.
For each block above, $r=q$ and $n=q$ or $q-1$.
Equation (CJ.11) gives $q-1\ge k$ for $k\ge3$,
and $q/2\le q\le2(q-1)$ for these $q$. Both norm factors
are therefore at most $C_h^{\rm win}q$.
Combining (CJ.11), (CJ.14) and (CJ.18), with
$D_h=\delta/(2C_h^{\rm win})>0$, proves for each block
\[
 \boxed{\log(V_i/V_j)\ge2q\log(D_h k)
       +\mathcal P_{i,j}
       +\sum_{N=i}^{j-1}\log(1+\phi_N^2/L^2).}
 \tag{CJ.19}
\]
Adding the two bounds and dividing by $q\log k$ proves
the four-volume lower limit at least four. This includes the
complete finite penalties before taking that consequence.

\subsection{The full relation cost and its retained zero directions}

For any admitted $i<j$, set $F=[b_{i+1},\ldots,b_j]$ and
$\Omega=\operatorname{diag}(\omega_{i+1},\ldots,\omega_j)$.
The exact update is $K_j-K_i=F\Omega^{-1}F^*$, so
\[
 V_i/V_j=\det(I+F\Omega^{-1}F^*G_i)
         =\det(I+\Omega^{-1}F^*G_iF).
 \tag{CJ.20}
\]
The rectangular determinant equality follows by eliminating either
off-diagonal block of the original block matrix
\[
 \begin{pmatrix}I&F\\-\Omega^{-1}F^*G_i&I\end{pmatrix}.
\]
On the new-column source $\C^{j-i}$ with metric $\Omega$, the
second cost is positive self-adjoint. Its eigenvalues are
nonnegative, including their zero multiplicities.
Write $L_{i,j}:\C^{i+1}\to\C^{j+1}$ for padding in the
same monic orthogonal basis, and
$T_{\rm new}:\C^{j-i}\to\C^{j+1}$ for the coordinate inclusion
into positions $i+1,\ldots,j$. The original relation map from
$(\C^{j-i},\Omega)$ to $(\mathcal P_j,O_j)$ is
$\mathfrak b=T_{\rm new}-L_{i,j}R_iF$.
The new columns and the lower source are orthogonal, and
$R_i^*O_iR_i=G_i$, so
\[
 B_j\mathfrak b=0,\qquad
 \mathfrak b^*O_j\mathfrak b=\Omega+F^*G_iF.
 \tag{CJ.21}
\]
Set $\mathcal D_N=\ker B_N=\chi\mathcal P_{N-q}$, with the
zero space when $N=q-1$. The new-coordinate component of $\mathfrak b$ is the identity;
it is therefore injective. Its image lies in
$\mathcal D_j\cap(L\mathcal D_i)^\perp$, whose dimension is
$j-i$, by the dimensions of the two nested relation spaces.
Thus it is an isomorphism onto that entire relation layer.
Reduction sends its labelled vectors to their receiving supported
zero; the lifted reduction fixes external absence $\tau$.

At zero tilt the full quartet gives an even original density.
Its $u$-polynomial $\psi=i^{-q}\chi(c+iu)$ satisfies
$\psi(-u)=(-1)^q\psi(u)$ because its complete root multiset
is stable under negation. Uniqueness of the monic orthogonal
polynomial gives $Q_n(-u)=(-1)^nQ_n(u)$: substituting $-u$
preserves all original integral pairings and the adjusted leading
coefficient. Monic division modulo $\psi$ preserves parity.
When $k$ is even, $q$ in (CJ.11) is odd. In the first block
the $q$ new polynomials $Q_q,\ldots,Q_{2q-1}$ contain
$(q+1)/2$ odd polynomials, while the odd quotient subspace has
dimension $(q-1)/2$. Their reduction columns are dependent.
The original $S=c+iu$ basis map and the nonzero column phases
$i^n$ preserve rank. Thus at least one of the $q$ cost eigenvalues
in that block is zero. It remains in the full determinant as a
factor one; (CJ.19)--(CJ.21) remain valid with every original
dimension and nilpotent direction retained.

For either length-$q$ block let $\lambda_a\ge0$ be all its cost
eigenvalues. Equations (CJ.19)--(CJ.20) imply
\[
 \prod_{a=1}^q(1+\lambda_a)\ge(D_hk)^{2q},\qquad
 \sum_{a=1}^q(1+\lambda_a)^s\ge q(D_hk)^{2s}
 \quad(s>0).
 \tag{CJ.22}
\]
The second assertion is the arithmetic-geometric mean inequality
for the $q$ positive numbers $(1+\lambda_a)^s$; it also obeys
the trivial lower bound $q$. Here is its exact restriction map.
Give the same coefficient space $E$ its two original metrics and
let $I_{i,j}:(E,G_i)\to(E,G_j)$ be the identity coefficient map.
Its Hilbert adjoint is
$T_{i,j}:(E,G_j)\to(E,G_i)$ with matrix $G_i^{-1}G_j$,
because $x^*G_jy=x^*G_i(G_i^{-1}G_j)y$ for every $x,y$.
Thus $Q=T_{i,j}I_{i,j}$ acts on $(E,G_i)$ and satisfies
\[
 Q^{-1}=G_j^{-1}G_i=K_jG_i
      =I+F\Omega^{-1}F^*G_i.
\]
It is positive self-adjoint in that original metric; its
quadratic form is $x^*G_iK_jG_ix>0$ for $x\ne0$.
Consequently its eigenvalues are exactly the $1+\lambda_a$
in (CJ.22), with all zero cost directions giving eigenvalue one.
Hence (CJ.22) is also an inverse-restriction trace
bound in its specified original metric. No individual positive
cost lower bound or strict-positive rank claim is inferred.

\subsection{Provenance and verification scope}

The complete cumulative sources CV and TVB prove the radius and
phase map used here, and the complete arithmetic source and AU
prove the actual density envelope (CJ.15). PR~25 supplies the
comparable-window constants and first-degree propagation. The
full first-window phase product (CJ.9), its finite penalties
(CJ.13)--(CJ.19), and their connection to the retained relation
costs are written out here. The owner review at the cited head
requires a nonnegative-eigenvalue correction and workflow-trigger
corrections; this chapter preserves that status and the complete
proposed parity proof. It performs no local Lean execution and
makes no new formal-verification claim. Its finite algebra and
analytic envelope calculations belong to this written proof.

\endgroup
